\documentclass[11pt]{article}
\usepackage[margin=1in]{geometry}
\usepackage[T1]{fontenc}
\usepackage{lmodern}
\usepackage{microtype}
\usepackage{amsmath,amssymb,amsthm,mathtools}
\usepackage{graphicx}
\usepackage{xcolor}
\usepackage{enumitem}
\usepackage[hidelinks]{hyperref}

\definecolor{statusblue}{RGB}{25,72,122}
\newtheorem{theorem}{Theorem}
\newtheorem{lemma}[theorem]{Lemma}
\newtheorem{proposition}[theorem]{Proposition}
\newtheorem{remark}[theorem]{Remark}
\newcommand{\E}{\mathcal E}
\newcommand{\M}{\mathsf M}
\newcommand{\N}{\mathbb N}
\newcommand{\Z}{\mathbb Z}
\newcommand{\Q}{\mathbb Q}

\title{Repeated generators force continuum many\\rapidly growing AF algebras}
\author{Candidate full solution packet for arXiv:2606.12866}
\date{9 August 2026}

\begin{document}
\maketitle

\begin{center}
\fcolorbox{statusblue}{statusblue!6}{\parbox{0.90\linewidth}{
\textbf{Status.} Candidate full solution, likely valid. The theorem below answers Problem~8.2 of Aguilar--Garcia--Knight--Marple--Spielberg affirmatively, subject to human review.
}}
\end{center}

\section{Source and open problem}

The source is K.~Aguilar, S.~R.~Garcia, E.~Knight, C.~Marple, and J.~Spielberg, \emph{Rapidly growing AF algebras}, arXiv:2606.12866 (2026)~\cite{source}. Problem~8.2 appears on PDF page~31:

\begin{center}
\includegraphics[width=0.97\linewidth]{figures/open_problem_crop.png}
\end{center}

In the source notation, let $k\ge 2$, let
\[
  \mathbf n=(n_1,\ldots,n_k)^T\in\N^k,
  \qquad 1=g_1<g_2<\cdots,
  \qquad \gamma_r=\frac{g_{r+1}}{g_r},
\]
where $g_r\mid g_{r+1}$ and
\[
  \sum_{r=1}^{\infty}\gamma_r^{-k}<\infty.
\]
At level $r$ the finite-dimensional algebra is
\[
  A_r=\bigoplus_{i=1}^k \M_{g_rn_i}.
\]
The ensemble $\E(\mathbf n,\mathbf g)$ consists of the chains whose standard unital connecting maps have nonnegative integral multiplicity matrices $X_r$ satisfying
\begin{equation}\label{eq:ensemble}
  X_r\mathbf n=\gamma_r\mathbf n.
\end{equation}
The question is whether a repeated entry of $\mathbf n$ forces uncountably many isomorphism classes among the limits.

\section{Proof intuition}

Two equal generators already support the variable $2\times2$ construction used by source Theorem~7.1. The other generators need not interact with them: put the variable block on the repeated pair and use diagonal multiplicity $\gamma_r$ on every remaining coordinate. This produces a direct sum of one variable simple AF algebra and finitely many fixed UHF algebras.

There are two points to verify. First, the $2\times2$ parameters must have the same parity as $\gamma_r$. A sparse sequence of distinct odd primes provides continuum many choices while respecting parity for both even and odd $\gamma_r$. Second, adjoining the same fixed UHF summands cannot identify two variable algebras. Indeed, an isomorphism of finite direct sums of simple algebras permutes their minimal nonzero ideals, so the common finite multiset of UHF summands cancels.

\section{A parity-safe continuum of two-generator limits}

The source uses, for $0<q_r<\gamma_r$ with $q_r\equiv\gamma_r\pmod 2$,
\begin{equation}\label{eq:block}
  p_r=\frac{\gamma_r+q_r}{2},
  \qquad
  Y_r(q)=
  \begin{pmatrix}
    p_r&\gamma_r-p_r\\
    \gamma_r-p_r&p_r
  \end{pmatrix}.
\end{equation}
Then $Y_r(q)(1,1)^T=\gamma_r(1,1)^T$. For a sequence $q=(q_r)$, write
\[
  a(q)=\prod_r q_r,
  \qquad G=\prod_r\gamma_r
\]
as generalized integers. Source Proposition~7.6 implies that if the abstract $K_0$-groups $H(q)$ and $H(q')$ of two such limits are isomorphic, then either
\begin{equation}\label{eq:source-invariant}
  a(q)\sim a(q')
  \quad\text{or}\quad
  a(q)\sim G\sim a(q'),
\end{equation}
where $\sim$ is equivalence of generalized integers up to finitely many finite prime exponents~\cite[Prop.~7.6]{source}.

\begin{lemma}[Parity-safe coding]\label{lem:coding}
For the fixed sequence $(\gamma_r)$ there is a family $\{q^S:S\in\mathfrak S\}$ of cardinality $2^{\aleph_0}$ such that:
\begin{enumerate}[label=(\roman*)]
\item $0<q_r^S<\gamma_r$ and $q_r^S\equiv\gamma_r\pmod 2$ for every sufficiently large $r$;
\item the generalized integers $a_S=\prod_rq_r^S$ are pairwise inequivalent;
\item no $a_S$ is equivalent to $G=\prod_r\gamma_r$.
\end{enumerate}
Consequently, the associated two-generator AF limits $B_S$ have pairwise non-isomorphic $K_0$-groups, and hence are pairwise non-isomorphic.
\end{lemma}

\begin{proof}
The summability hypothesis gives $\gamma_r\to\infty$. Discard a finite initial segment so that $\gamma_r>2$ thereafter; this does not change an inductive limit. Choose increasing indices $r_j$ and pairwise distinct odd primes $\ell_j$ so that
\[
  2\ell_j<\gamma_{r_j}.
\]
For indices in the retained tail, put
\[
  c_r=\begin{cases}
       1,&\gamma_r\text{ odd},\\
       2,&\gamma_r\text{ even}.
      \end{cases}
\]
For $S\subseteq\N$, define
\[
 q_r^S=
 \begin{cases}
   c_{r_j}\ell_j,&r=r_j\text{ and }j\in S,\\
   c_r,&\text{otherwise}.
 \end{cases}
\]
These choices have the required parity and lie strictly between $0$ and $\gamma_r$. The exponent of the odd prime $\ell_j$ in $a_S$ equals $1$ exactly when $j\in S$. Hence
\[
  a_S\sim a_T \quad\Longleftrightarrow\quad S\mathbin{\triangle}T\text{ is finite}.
\]
The quotient $\mathcal P(\N)/\mathrm{Fin}$ has cardinality $2^{\aleph_0}$. Choose one representative from each equivalence class. At most one of the resulting classes can be equivalent to the fixed generalized integer $G$; discard it. This proves (i)--(iii).

For the finitely many deleted indices use the diagonal block $\gamma_rI_2$; equivalently, begin the system at the retained tail. Every later matrix in~\eqref{eq:block} has all entries strictly positive, so $B_S$ is simple. Finally, source Proposition~7.6 and (iii) rule out the cross-alternative in~\eqref{eq:source-invariant}, while (ii) rules out the first alternative. Thus $K_0(B_S)\not\cong K_0(B_T)$ for $S\ne T$.
\end{proof}

\begin{remark}[Repair of a parity omission]
The printed proof of source Theorem~7.1 selects $q_r$ to be prime. When $\gamma_r$ is even, this does not ensure the necessary congruence $q_r\equiv\gamma_r\pmod 2$; an odd prime makes $p_r=(\gamma_r+q_r)/2$ nonintegral. Lemma~\ref{lem:coding} supplies a parity-safe replacement and, in particular, a corrected continuum argument for the two-generator theorem used here.
\end{remark}

\section{Full resolution of Problem 8.2}

\begin{theorem}\label{thm:main}
Suppose that the generating vector $\mathbf n$ has a repeated entry. Then
\[
  \E(\mathbf n,\mathbf g)
\]
contains $2^{\aleph_0}$ pairwise non-isomorphic AF algebras. Thus Problem~8.2 of~\cite{source} has an affirmative answer.
\end{theorem}

\begin{proof}
After permuting coordinates, assume $n_1=n_2=n$. Let $\{B_S:S\in\mathfrak S\}$ be the family from Lemma~\ref{lem:coding}. At each sufficiently late level define the $k\times k$ multiplicity matrix
\begin{equation}\label{eq:bigblock}
  X_r^S=operatorname{diag}\bigl(Y_r(q^S),\gamma_r,\ldots,\gamma_r\bigr).
\end{equation}
At the finitely many earlier levels use $X_r^S=\gamma_r I_k$. Because $n_1=n_2$, the first two rows of~\eqref{eq:ensemble} follow from
\[
  Y_r(q^S)(n,n)^T=\gamma_r(n,n)^T,
\]
and each remaining row is immediate. Thus $X_r^S\mathbf n=\gamma_r\mathbf n$, and every constructed chain belongs to the original ensemble.

All connecting maps are block diagonal with respect to
\[
  \bigl(\M_{g_rn}\oplus\M_{g_rn}\bigr)
  \oplus\bigoplus_{j=3}^k\M_{g_rn_j}.
\]
Taking the direct limit coordinatewise gives
\begin{equation}\label{eq:decomp}
  A_S\cong B_S\oplus U_3\oplus\cdots\oplus U_k,
\end{equation}
where $U_j$ is the fixed UHF limit of
\[
  \M_{g_rn_j}\xrightarrow{\,\gamma_r\,}\M_{g_{r+1}n_j}.
\]
The algebra $B_S$ is simple because all sufficiently late entries of its Bratteli matrices are positive, and every $U_j$ is simple.

Suppose $A_S\cong A_T$. In a finite direct sum of simple $C^*$-algebras, the minimal nonzero closed two-sided ideals are exactly the summands. An isomorphism therefore induces equality of the finite multisets of isomorphism classes of simple summands in~\eqref{eq:decomp}. Cancelling the common multiset
\[
  \{[U_3],\ldots,[U_k]\}
\]
gives $[B_S]=[B_T]$. This contradicts Lemma~\ref{lem:coding} whenever $S\ne T$. Hence the $A_S$ are pairwise non-isomorphic, and $|\mathfrak S|=2^{\aleph_0}$.
\end{proof}

\section{Verification notes and limitations}

\begin{enumerate}[leftmargin=*,itemsep=4pt]
\item \textbf{Membership in the ensemble.} Equation~\eqref{eq:bigblock} has nonnegative integral entries, and the row equations $X_r^S\mathbf n=\gamma_r\mathbf n$ were checked explicitly.
\item \textbf{Parity and endpoints.} The coding uses $q_r=1$ or an odd prime when $\gamma_r$ is odd, and $q_r=2$ or twice an odd prime when $\gamma_r$ is even. The strict inequality follows from $2\ell_j<\gamma_{r_j}$ and from deleting the finite set with $\gamma_r\le2$.
\item \textbf{Continuum cardinality.} The invariant records subsets of $\N$ modulo finite symmetric difference, not merely individual binary sequences. This quotient still has cardinality $2^{\aleph_0}$.
\item \textbf{K-theory input.} The only non-elementary imported step is source Proposition~7.6. The proof deliberately removes the possible class $a_S\sim G$, so the proposition's coordinate-swap alternative cannot collapse two retained members.
\item \textbf{Cancellation.} No cancellation theorem for arbitrary $C^*$-algebras is assumed. It is the elementary uniqueness of the finite multiset of minimal ideals in a finite direct sum of simple algebras.
\item \textbf{Scope.} The theorem fully resolves source Problem~8.2. It does not answer Problem~8.1 for generating vectors with no repeated entries, nor the source's K-theory, simplicity-probability, or UHF-embedding problems.
\end{enumerate}

No numerical computation is part of the proof. A bounded novelty search on 9 August 2026 covered the run registry and packet indexes, the locally parsed arXiv corpus, and web queries using the exact title, arXiv id, Problem~8.2 wording, ``repeated generator'', and ``uncountably many non-isomorphic AF algebras''. These searches found the source paper but no later answer or matching theorem. Since the source is recent, the novelty assessment is necessarily provisional.

\section{Human-review recommendation}

The proof appears complete and should be reviewed as a candidate full solution. The two highest-value checks are the use of source Proposition~7.6 after deletion of the cross-equivalence class and the minimal-ideal cancellation in~\eqref{eq:decomp}. A reviewer should also note that Lemma~\ref{lem:coding} repairs the parity omission in the source's proof of Theorem~7.1.

\begin{thebibliography}{9}
\bibitem{source}
K.~Aguilar, S.~R.~Garcia, E.~Knight, C.~Marple, and J.~Spielberg,
\emph{Rapidly growing AF algebras}, arXiv:2606.12866 (2026), especially Sections~3 and~7, Proposition~7.6, and Problem~8.2.
\url{https://arxiv.org/abs/2606.12866}.
\end{thebibliography}

\end{document}
